\begin{solapartat}
\figura[0.45]{sol_fig1.png}

El camp magnètic que crea el fil infinit en el centre de l'espira és: $B_{\mathit{fil}\,\infty} = \dfrac{\mu_0 I}{2\pi d}$

\punts{0,3} $\left|B_{\mathit{fil}\,\infty}\right| = \dfrac{4\pi\times 10^{-7}\times 2}{2\pi\times 0,05} = 8,0\times 10^{-6}\ T$

\punts{0,2} i entra al paper $\otimes$

En notació vectorial: $\vec{B}_{\mathit{fil}\,\infty} = -8,0\times 10^{-6}\,\vec{k}\ \ T$

El camp magnètic que crea l'espira al centre és: $B_{\mathit{espira}} = \dfrac{\mu_0 I}{2R}$

\punts{0,3} $\left|B_{\mathit{espira}}\right| = \dfrac{4\pi\times 10^{-7}\times 0,5}{2\times 1,0\times 10^{-2}} = 3,1\times 10^{-5}\ T$ i surt del paper $\odot$ \punts{0,2}

En notació vectorial: $\vec{B}_{\mathit{espira}} = 3,1\times 10^{-5}\,\vec{k}\ \ T$
\end{solapartat}

\begin{solapartat}
El camp magnètic total és la suma vectorial del camp creat pel fil infinit i l'espira:

En notació vectorial:\\
% DUBTE: com a suma vectorial hauria de ser B_espira + B_fil (amb B_fil = -8,0·10^-6 k); la pauta escriu la resta amb els mòduls. Es transcriu tal com és.
$\vec{B}_{\mathit{total}} = \vec{B}_{\mathit{espira}} - \vec{B}_{\mathit{fil}\,\infty} = 3,1\times 10^{-5}\,\vec{k} - 8,0\times 10^{-6}\,\vec{k} = 2,3\times 10^{-5}\,\vec{k}\ \ T$

\punts{0,2} Mòdul: $\left|B_{\mathit{total}}\right| = 2,3\times 10^{-5}\ T$ i surt del paper $\odot$ \punts{0,2}

\punts{0,2} Si $B = 0 \Rightarrow B'_{\mathit{espira}} = B_{\mathit{fil}} = 8,0\times 10^{-6}\ T$

\punts{0,4} $B'_{\mathit{espira}} = \dfrac{\mu_0 I'}{2R} \Rightarrow I' = \dfrac{2RB'_{\mathit{espira}}}{\mu_0} = \dfrac{2\times 0,01\times 8,0\times 10^{-6}}{4\pi\times 10^{-7}} = 0,13\ A$
\end{solapartat}
